
\documentclass[11pt,a4paper]{article}

\usepackage[margin=0.8in]{geometry}
\usepackage[T1]{fontenc}
\usepackage{lmodern}
\usepackage{xcolor}
\usepackage{hyperref}
\usepackage{enumitem}
\usepackage{titlesec}
\usepackage{tabularx}
\usepackage{fancyhdr}
\usepackage{tcolorbox}
\usepackage{amssymb}

\definecolor{darkblue}{RGB}{30,65,110}
\definecolor{lightblue}{RGB}{235,243,252}
\definecolor{lightgray}{RGB}{245,245,245}
\definecolor{darkgreen}{RGB}{35,105,70}

\hypersetup{colorlinks=true,linkcolor=darkblue,urlcolor=darkblue}

\titleformat{\section}{\Large\bfseries\color{darkblue}}{\thesection}{0.6em}{}
\titleformat{\subsection}{\large\bfseries\color{darkblue}}{\thesubsection}{0.5em}{}
\titleformat{\subsubsection}{\normalsize\bfseries\color{darkgreen}}{\thesubsubsection}{0.5em}{}

\pagestyle{fancy}
\fancyhf{}
\lhead{Unit 1 -- Object-Oriented Thinking}
\rhead{Course Tracker}
\cfoot{\thepage}

\newtcolorbox{conceptbox}{
    colback=lightblue,colframe=darkblue,boxrule=0.7pt,arc=2mm,
    left=2mm,right=2mm,top=1.5mm,bottom=1.5mm
}
\newtcolorbox{exambox}{
    colback=lightgray,colframe=black!40,boxrule=0.5pt,arc=2mm,
    left=2mm,right=2mm,top=1.5mm,bottom=1.5mm
}

\newcommand{\checkitem}{$\square$}

\title{\vspace{-1cm}\textbf{Unit 1 -- Object-Oriented Thinking}\\
\large Complete Topic \& Subtopic Tracker}
\author{Course Notes}
\date{}

\begin{document}
\maketitle

\begin{conceptbox}
\textbf{Syllabus scope:} This document covers the Object-Oriented Thinking
portion of Unit 1: encapsulation, abstraction, inheritance vs composition,
and polymorphism.
\end{conceptbox}

\tableofcontents
\newpage

% =========================================================
\section{Object-Oriented Thinking}

\subsection{Study Tracker}

\begin{itemize}[label=\checkitem]
    \item Understand the object-oriented way of modeling a problem
    \item Identify objects and their responsibilities
    \item Understand state and behavior
    \item Understand classes and objects
    \item Understand encapsulation
    \item Understand abstraction
    \item Understand inheritance
    \item Compare inheritance and composition
    \item Understand polymorphism
\end{itemize}

\subsection{What is Object-Oriented Thinking?}

Object-oriented thinking models a system using objects that represent
entities or concepts in the problem domain.

An object generally combines:

\begin{itemize}
    \item \textbf{State} -- the data the object contains.
    \item \textbf{Behavior} -- the operations the object can perform.
\end{itemize}

For example, a \texttt{BankAccount} could have:

\begin{verbatim}
State:
    accountNumber
    balance

Behavior:
    deposit()
    withdraw()
    getBalance()
\end{verbatim}

The goal is to decide which objects should exist and which responsibilities
belong to each object.

\subsection{Class vs Object}

A \textbf{class} describes the structure and behavior of a type.

An \textbf{object} is an instance of that class.

Conceptually:

\[
\text{Class} \rightarrow \text{Object}
\]

For example:

\begin{verbatim}
Class:   Car
Objects: car1, car2, car3
\end{verbatim}

A class defines what an object can contain and do; an object represents a
particular instance.

% =========================================================
\section{Encapsulation}

\subsection{Study Tracker}

\begin{itemize}[label=\checkitem]
    \item Define encapsulation
    \item Understand data hiding
    \item Understand access control
    \item Understand public/private/protected access
    \item Understand why behavior should control state changes
    \item Identify examples of encapsulation
\end{itemize}

\subsection{Definition}

Encapsulation means bundling an object's data and the operations that work
with that data together while controlling access to the internal state.

A simple conceptual model is:

\[
\boxed{\text{Object} = \text{State} + \text{Behavior}}
\]

with controlled access to the state.

\subsection{Data Hiding}

An object's internal data does not always need to be directly accessible.

For example, instead of allowing arbitrary changes to a bank balance, the
class can expose operations such as:

\begin{verbatim}
deposit()
withdraw()
getBalance()
\end{verbatim}

These operations can enforce rules.

\subsection{Access Control}

Typical object-oriented access levels include:

\begin{itemize}
    \item \textbf{Public}: accessible through the public interface.
    \item \textbf{Private}: accessible only within the class.
    \item \textbf{Protected}: accessible within the class and appropriate
    derived classes, depending on the language.
\end{itemize}

The exact syntax depends on the programming language.

\subsection{Why Encapsulation Matters}

Encapsulation can:

\begin{itemize}
    \item protect internal state,
    \item prevent invalid direct modifications,
    \item localize implementation details,
    \item provide a controlled interface,
    \item make future implementation changes easier.
\end{itemize}

\begin{exambox}
\textbf{Exam definition:} Encapsulation is the bundling of data and related
behavior together while controlling access to the object's internal state.
\end{exambox}

% =========================================================
\section{Abstraction}

\subsection{Study Tracker}

\begin{itemize}[label=\checkitem]
    \item Define abstraction
    \item Understand essential vs non-essential details
    \item Understand interfaces
    \item Understand abstract classes/interfaces conceptually
    \item Distinguish abstraction from implementation
    \item Compare abstraction with encapsulation
\end{itemize}

\subsection{Definition}

Abstraction means exposing the essential operations or concepts of an object
while hiding unnecessary implementation details.

The user of an abstraction focuses on:

\[
\boxed{\text{What can this object do?}}
\]

rather than:

\[
\boxed{\text{How is every internal step implemented?}}
\]

\subsection{Example}

A payment interface might expose:

\begin{verbatim}
pay(amount)
\end{verbatim}

The caller does not necessarily need to know whether the implementation uses
a particular banking API, database operation, or network protocol.

\subsection{Interface as an Abstraction}

An interface defines a contract that an implementation must satisfy.

Conceptually:

\begin{verbatim}
Payment
    |
    +-- CardPayment
    +-- UPIPayment
    +-- CashPayment
\end{verbatim}

The caller can depend on the common payment abstraction rather than every
concrete implementation.

\subsection{Abstraction vs Encapsulation}

\begin{tabularx}{\textwidth}{|p{0.25\textwidth}|X|}
\hline
\textbf{Abstraction} & Focuses on exposing essential behavior and hiding unnecessary implementation details.\\
\hline
\textbf{Encapsulation} & Focuses on bundling state and behavior and controlling access to internal state.\\
\hline
\end{tabularx}

\begin{exambox}
\textbf{Easy memory aid:}\\
Abstraction asks: \textit{``What should the outside world see?''}\\
Encapsulation asks: \textit{``How do I protect and control the inside?''}
\end{exambox}

% =========================================================
\section{Inheritance}

\subsection{Study Tracker}

\begin{itemize}[label=\checkitem]
    \item Define inheritance
    \item Understand base and derived classes
    \item Understand reuse through inheritance
    \item Understand method overriding conceptually
    \item Understand ``is-a'' relationships
    \item Understand appropriate and inappropriate inheritance
\end{itemize}

\subsection{Definition}

Inheritance allows one class to derive from another class.

Conceptually:

\[
\text{Base Class} \rightarrow \text{Derived Class}
\]

For example:

\[
Vehicle
\rightarrow
\begin{cases}
Car\\
Bike
\end{cases}
\]

The derived class can inherit appropriate state and behavior and can add or
specialize behavior.

\subsection{Is-A Relationship}

Inheritance should normally represent a meaningful ``is-a'' relationship.

Examples:

\[
Car \text{ is a } Vehicle
\]

\[
Dog \text{ is an } Animal
\]

A class hierarchy should preserve the meaning and expectations of the
base abstraction.

\subsection{Method Overriding}

A derived class may provide a specialized implementation of a behavior
defined by the base class.

Conceptually:

\begin{verbatim}
Animal
    makeSound()

Dog
    makeSound()  --> dog-specific behavior

Cat
    makeSound()  --> cat-specific behavior
\end{verbatim}

This becomes important when understanding runtime polymorphism.

\subsection{Advantages of Inheritance}

Potential benefits:

\begin{itemize}
    \item reuse of common behavior,
    \item representation of genuine type hierarchies,
    \item support for polymorphism,
    \item specialization of common abstractions.
\end{itemize}

\subsection{Risks of Inheritance}

Inheritance can create strong coupling between parent and child classes.

Problems can occur when:

\begin{itemize}
    \item the ``is-a'' relationship is artificial,
    \item the hierarchy becomes too deep,
    \item subclasses depend heavily on parent implementation,
    \item inherited behavior does not make sense for the subclass.
\end{itemize}

This connects directly to the Liskov Substitution Principle in Unit 2.

% =========================================================
\section{Inheritance vs Composition}

\subsection{Study Tracker}

\begin{itemize}[label=\checkitem]
    \item Define inheritance-based design
    \item Define composition-based design
    \item Understand ``is-a'' vs ``has-a''
    \item Compare flexibility
    \item Identify when composition is preferable
    \item Understand how this connects to design patterns and SOLID
\end{itemize}

\subsection{Inheritance: ``Is-A''}

Inheritance models a type hierarchy.

\[
Car \text{ is a } Vehicle
\]

The child derives from the parent.

\subsection{Composition: ``Has-A''}

Composition builds an object using other objects.

\[
Car \text{ has an } Engine
\]

Conceptually:

\begin{verbatim}
Car
 |
 +-- Engine
\end{verbatim}

The Car does not need to become a subclass of Engine.

\subsection{Comparison}

\begin{tabularx}{\textwidth}{|p{0.23\textwidth}|X|X|}
\hline
 & \textbf{Inheritance} & \textbf{Composition}\\
\hline
Relationship & is-a & has-a\\
\hline
Structure & parent-child hierarchy & object contains/uses another object\\
\hline
Coupling & often stronger & often more flexible\\
\hline
Reuse & through inherited behavior & through contained components\\
\hline
Typical example & Car is a Vehicle & Car has an Engine\\
\hline
\end{tabularx}

\subsection{Why Composition is Often Preferred}

Composition can allow a component to be replaced or changed independently.

For example, instead of creating many subclasses for different behaviors:

\begin{verbatim}
CarWithX
CarWithY
CarWithZ
\end{verbatim}

a design can sometimes use interchangeable components:

\begin{verbatim}
Car
 |
 +-- Engine
 +-- PaymentStrategy
 +-- NavigationSystem
\end{verbatim}

This can reduce the size of inheritance hierarchies.

\begin{exambox}
\textbf{Exam question: Why prefer composition over inheritance?}\\
A common answer is that composition can provide greater flexibility and
lower coupling because behavior can be assembled from independent objects
instead of being permanently tied to an inheritance hierarchy.
\end{exambox}

% =========================================================
\section{Polymorphism}

\subsection{Study Tracker}

\begin{itemize}[label=\checkitem]
    \item Define polymorphism
    \item Understand the common interface idea
    \item Understand method overriding
    \item Understand runtime polymorphism
    \item Understand compile-time polymorphism
    \item Connect polymorphism with abstraction and inheritance
\end{itemize}

\subsection{Definition}

Polymorphism means that the same interface or operation can represent
different behavior depending on the object involved.

The central idea is:

\[
\boxed{\text{One interface} \rightarrow \text{multiple implementations}}
\]

\subsection{Runtime Polymorphism}

Runtime polymorphism occurs when the implementation selected depends on the
actual object at runtime.

Conceptually:

\begin{verbatim}
Animal* animal = Dog;
animal->makeSound();
\end{verbatim}

The common interface is \texttt{Animal}, while the actual implementation
can be supplied by \texttt{Dog}.

Another object could be:

\begin{verbatim}
Animal* animal = Cat;
animal->makeSound();
\end{verbatim}

The same operation can therefore produce different behavior.

\subsection{Compile-Time Polymorphism}

Depending on the language, compile-time polymorphism can include mechanisms
such as function overloading or operator overloading.

Example concept:

\begin{verbatim}
add(int, int)
add(double, double)
\end{verbatim}

The operation has the same conceptual name but different parameter forms.

\subsection{Polymorphism and Abstraction}

Abstraction defines a common interface.

Polymorphism allows different implementations to be used through that common
interface.

\[
\text{Abstraction}
\rightarrow
\text{Common Interface}
\rightarrow
\text{Polymorphism}
\rightarrow
\text{Different Implementations}
\]

\subsection{Polymorphism and Inheritance}

Inheritance is one common way to implement runtime polymorphism.

However, the broader design idea is programming to a common abstraction and
allowing different implementations to satisfy it.

% =========================================================
\section{Connections Between the OOP Concepts}

\subsection{The Big Picture}

The five topics in this section work together:

\[
\boxed{
Encapsulation
+
Abstraction
+
Inheritance
+
Composition
+
Polymorphism
}
\]

A typical design may use several simultaneously.

For example, consider a payment system:

\begin{verbatim}
Payment
   |
   +-- CardPayment
   +-- UPIPayment
   +-- CashPayment
\end{verbatim}

The common \texttt{Payment} abstraction defines what payment implementations
must provide.

Encapsulation keeps implementation details inside each class.

Polymorphism allows the application to call the common payment operation
without knowing which concrete implementation is being used.

Composition can be used to build the payment system from other components.

% =========================================================
\section{Important Comparisons}

\subsection{Encapsulation vs Abstraction}

\begin{tabularx}{\textwidth}{|p{0.25\textwidth}|X|}
\hline
\textbf{Encapsulation} & Controls access to internal state and bundles state with related behavior.\\
\hline
\textbf{Abstraction} & Shows essential behavior while hiding unnecessary implementation details.\\
\hline
\end{tabularx}

\subsection{Inheritance vs Composition}

\begin{tabularx}{\textwidth}{|p{0.25\textwidth}|X|}
\hline
\textbf{Inheritance} & Models an ``is-a'' relationship through a class hierarchy.\\
\hline
\textbf{Composition} & Models a ``has-a'' relationship by combining objects.\\
\hline
\end{tabularx}

\subsection{Overloading vs Overriding}

\begin{tabularx}{\textwidth}{|p{0.25\textwidth}|X|}
\hline
\textbf{Overloading} & Multiple functions/methods with the same name but different parameter lists; commonly associated with compile-time polymorphism.\\
\hline
\textbf{Overriding} & A derived implementation replaces/specializes inherited behavior; commonly associated with runtime polymorphism.\\
\hline
\end{tabularx}

% =========================================================
\section{Exam-Oriented Questions}

\begin{enumerate}
    \item What is object-oriented thinking?
    \item What is the difference between a class and an object?
    \item Define encapsulation with an example.
    \item What is data hiding?
    \item Define abstraction.
    \item Differentiate abstraction and encapsulation.
    \item What is inheritance?
    \item Explain the ``is-a'' relationship.
    \item What is polymorphism?
    \item Differentiate compile-time and runtime polymorphism.
    \item What is method overriding?
    \item What is composition?
    \item Differentiate inheritance and composition.
    \item Why can composition provide more flexibility than inheritance?
    \item How are abstraction, inheritance, and polymorphism related?
    \item Give a real-world example where inheritance is appropriate.
    \item Give a real-world example where composition is appropriate.
    \item Identify the OOP concepts present in a given class design.
\end{enumerate}

% =========================================================
\section{Final Revision Checklist}

\subsection*{Object-Oriented Thinking}

\begin{itemize}[label=\checkitem]
    \item I can explain object-oriented thinking.
    \item I understand classes and objects.
    \item I understand state and behavior.
\end{itemize}

\subsection*{Encapsulation}

\begin{itemize}[label=\checkitem]
    \item I can define encapsulation.
    \item I understand data hiding.
    \item I understand access control.
    \item I can give a real-world example.
\end{itemize}

\subsection*{Abstraction}

\begin{itemize}[label=\checkitem]
    \item I can define abstraction.
    \item I understand interfaces conceptually.
    \item I can distinguish abstraction from encapsulation.
\end{itemize}

\subsection*{Inheritance}

\begin{itemize}[label=\checkitem]
    \item I can define inheritance.
    \item I understand base and derived classes.
    \item I understand the ``is-a'' relationship.
    \item I understand method overriding.
    \item I know the risks of inappropriate inheritance.
\end{itemize}

\subsection*{Composition}

\begin{itemize}[label=\checkitem]
    \item I can define composition.
    \item I understand the ``has-a'' relationship.
    \item I can compare composition with inheritance.
    \item I understand why composition can improve flexibility.
\end{itemize}

\subsection*{Polymorphism}

\begin{itemize}[label=\checkitem]
    \item I can define polymorphism.
    \item I understand runtime polymorphism.
    \item I understand compile-time polymorphism.
    \item I understand overriding vs overloading.
    \item I can explain polymorphism using a real-world example.
\end{itemize}

\begin{conceptbox}
\textbf{Final mental model:}

\[
\boxed{
\text{Encapsulation}
\rightarrow
\text{protect the inside}
}
\]

\[
\boxed{
\text{Abstraction}
\rightarrow
\text{show what matters}
}
\]

\[
\boxed{
\text{Inheritance}
\rightarrow
\text{is-a relationship}
}
\]

\[
\boxed{
\text{Composition}
\rightarrow
\text{has-a relationship}
}
\]

\[
\boxed{
\text{Polymorphism}
\rightarrow
\text{one interface, different behavior}
}
\]
\end{conceptbox}

\end{document}
